% Exercise 2.3-7
We must decide whether \emph{two} elements of $S$ have a sum equal to $x$, not
whether any single element equals $x$. We first sort $S$ with
\textsc{Merge-Sort}, and then sweep the sorted array inwards from both ends.

\begin{algorithmic}[1]
  \TITLE{\textsc{Sum-Exists}$(S, x)$}
  \STATE let $A[1 \ltwodots n]$ be the elements of $S$
  \STATE $\textsc{Merge-Sort}(A, 1, n)$
  \STATE $\mathit{low} = 1$
  \STATE $\mathit{high} = n$
  \WHILE{$\mathit{low} < \mathit{high}$}
  \STATE $\mathit{sum} = A[\mathit{low}] + A[\mathit{high}]$
  \IF{$\mathit{sum} == x$}
  \RETURN \textrm{TRUE}
  \ELSIF{$\mathit{sum} < x$}
  \STATE $\mathit{low} = \mathit{low} + 1$
  \ELSE
  \STATE $\mathit{high} = \mathit{high} - 1$
  \ENDIF
  \ENDWHILE

  \RETURN \textrm{FALSE}
\end{algorithmic}

The condition $\mathit{low} < \mathit{high}$ guarantees that the two elements
summed are distinct positions of $A$. If $A[\mathit{low}] + A[\mathit{high}] <
x$ then no element paired with $A[\mathit{low}]$ can reach $x$, since
$A[\mathit{high}]$ is the largest remaining candidate, and so $\mathit{low}$ may
safely be advanced; the symmetric argument justifies decrementing
$\mathit{high}$. Every iteration therefore discards an element that cannot
participate in any solution, and no pair summing to $x$ is ever lost.

\textsc{Merge-Sort} takes $\Theta(n\lg{n})$ time and the \textbf{while} loop of
lines 5-13 performs at most $n$ iterations of constant cost, so the total
running time is
\[
  \Theta(n\lg{n}) + \Theta(n) = \Theta(n\lg{n}).
\]
